\documentclass[../../main.tex]{subfiles}

\begin{document}

\chapter{Memulai Assembly: TASM, TLINK, COM, EXE, dan Direktif}
\label{chap:memulai-assembly}

% =========================================================
% OUTCOME PEMBELAJARAN KHUSUS BAB
% =========================================================

\begin{subcpmkbox}
\vspace{1.1em}
\SubCPMK{2.1} --- Menjelaskan struktur program assembly
(\texttt{.MODEL}, \texttt{ORG 100h}, \texttt{.CODE}, \texttt{.DATA},
label, komentar) dan alur translasi \texttt{.ASM} $\rightarrow$
\texttt{.OBJ} $\rightarrow$ \texttt{.EXE}/\texttt{.COM}.

\vspace{0.4em}
\textbf{CPMK yang didukung:} \CPMK{02} \quad\textbullet\quad
\textbf{CPL yang didukung:} \CPL{02}
\end{subcpmkbox}

\begin{tujuanbox}
\vspace{1.1em}
Setelah mempelajari bab ini, mahasiswa mampu:
\begin{enumerate}
    \item menjelaskan rantai proses dari berkas \file{.ASM} sampai
          program dapat dijalankan;
    \item menjelaskan peran editor, TASM, TLINK, dan \eng{loader};
    \item membedakan program \file{.COM} dan \file{.EXE};
    \item menjelaskan fungsi direktif \dir{.MODEL}, \dir{ORG 100h},
          \dir{.CODE}, \dir{.DATA}, dan \dir{.STACK};
    \item menulis label dan komentar sesuai aturan assembler;
    \item menggunakan direktif data \dir{DB}, \dir{DW}, \dir{DD},
          \dir{EQU}, dan \dir{DUP};
    \item menulis, meng-assemble, men-link, dan menjalankan program
          assembly sederhana.
\end{enumerate}
\end{tujuanbox}

\begin{materibox}
\vspace{1.1em}
\begin{enumerate}
    \item Rantai proses: editor, TASM, TLINK, dan \eng{loader}.
    \item Perbedaan program \file{.COM} dan \file{.EXE}.
    \item Struktur program: \dir{.MODEL}, \dir{.STACK}, \dir{.DATA},
          \dir{.CODE}, label, komentar, dan \dir{END}.
    \item \dir{ORG 100h} dan \eng{Program Segment Prefix}.
    \item Direktif data: \dir{DB}, \dir{DW}, \dir{DD}, \dir{EQU},
          dan \dir{DUP}.
    \item Alur translasi dan pesan kesalahan assembler.
\end{enumerate}
\end{materibox}

\subfile{section-01}
\subfile{section-02}
\subfile{section-03}

\section{Contoh Program}
\label{sec:contoh-bab7}

Dua contoh berikut memperlihatkan cara kerja format instruksi pada
program nyata: membongkar program dengan \eng{disassembler}, lalu
membandingkan jumlah byte dua bentuk penulisan yang setara.

\subsection{Contoh 1: Membongkar Program dengan Perintah \code{U}}

Program pada Kode~\ref{lst:opcode-demo} sengaja memuat lima
instruksi dengan bentuk operand berbeda, sehingga hasil bongkarnya
dapat dibandingkan dengan kode sumbernya baris demi baris.

\lstinputlisting[
    language=MASM,
    caption={Program \file{opcode-demo.asm} untuk latihan
    \eng{disassembly}.},
    label={lst:opcode-demo}
]{\CodeDir/opcode-demo.asm}

Setelah program dirakit dan dijalankan sekali, program tersebut
dibongkar dengan \file{DEBUG}. Perintah \code{U} dimulai pada alamat
titik masuk program. Hasilnya ditampilkan pada
Kode~\ref{lst:bongkar-opcode-demo}.

\begin{session}[caption={Hasil perintah \code{U} pada
\file{opcode-demo.exe}.}, label={lst:bongkar-opcode-demo}]
-u
0B2A:0000 B82A0B        MOV     AX,0B2A
0B2A:0003 8ED8          MOV     DS,AX
0B2A:0005 B80500        MOV     AX,0005
0B2A:0008 BB0100        MOV     BX,0001
0B2A:000B 03C3          ADD     AX,BX
0B2A:000D 89C1          MOV     CX,AX
0B2A:000F 2D0200        SUB     AX,0002
0B2A:0012 B409          MOV     AH,09
0B2A:0014 8D160000      LEA     DX,[0000]
0B2A:0018 CD21          INT     21
0B2A:001A B44C          MOV     AH,4C
0B2A:001C B000          MOV     AL,00
0B2A:001E CD21          INT     21
\end{session}

Perbandingan kode sumber dan hasil bongkar tersebut dirangkum pada
Tabel~\ref{tab:banding-sumber-bongkar}.

\begin{table}[H]
    \centering
    \caption{Perbandingan kode sumber dan hasil \eng{disassembly}
    pada Kode~\ref{lst:bongkar-opcode-demo}.}
    \label{tab:banding-sumber-bongkar}
    \small
    \begin{tabular}{@{}p{3.6cm}p{2.2cm}cp{4.0cm}@{}}
        \toprule
        \textbf{Kode sumber} & \textbf{Kode mesin} &
        \textbf{Byte} & \textbf{Field yang tampak} \\
        \midrule
        \lstasm|MOV AX, 0005h| & \hex{B8 05 00} & 3 &
        Opcode penambah register, immediate 16 bit \\
        \addlinespace
        \lstasm|MOV BX, 0001h| & \hex{BB 01 00} & 3 &
        Nomor register \reg{BX} = \dec{3}, jadi \hex{B8} $+$ \dec{3} \\
        \addlinespace
        \lstasm|ADD AX, BX| & \hex{03 C3} & 2 &
        ModR/M \bin{11 000 011}: \code{reg} = \reg{AX},
        \code{r/m} = \reg{BX} \\
        \addlinespace
        \lstasm|MOV CX, AX| & \hex{89 C1} & 2 &
        Opcode \hex{89} arah \code{r/m} $\leftarrow$ \code{reg};
        \code{reg} = \reg{AX}, \code{r/m} = \reg{CX} \\
        \addlinespace
        \lstasm|SUB AX, 0002h| & \hex{2D 02 00} & 3 &
        Opcode khusus \reg{AX} dengan immediate 16 bit \\
        \addlinespace
        \lstasm|LEA DX, pesan| & \hex{8D 16 xx xx} & 4 &
        ModR/M \hex{16}: \code{mod} = \bin{00},
        \code{r/m} = \bin{110} (alamat langsung 16 bit) \\
        \bottomrule
    \end{tabular}
\end{table}

\begin{infobox}
Instruksi \lstasm|LEA DX, pesan| tertulis sebagai
\lstasm|LEA DX,[0000]|. Disassembler menampilkan \eng{displacement}
apa adanya, tanpa mengetahui nama simbol \code{pesan}. Nama simbol
hanya hidup di tabel simbol assembler dan hilang seluruhnya pada kode
mesin (Bab~\ref{chap:translasi}).
\end{infobox}

\subsection{Contoh 2: Dua Bentuk Penulisan, Ukuran Berbeda}

Contoh kedua memperlihatkan bahwa dua penulisan yang menghasilkan
keadaan register sama dapat menempati jumlah byte berbeda. Program
pada Kode~\ref{lst:ukuran-instruksi} memuat kedua bentuk tersebut
secara berdampingan dalam satu program \file{.COM}.

\lstinputlisting[
    language=MASM,
    caption={Program \file{ukuran-instruksi.asm} membandingkan dua
    bentuk penulisan.},
    label={lst:ukuran-instruksi}
]{\CodeDir/ukuran-instruksi.asm}

Hasil bongkar kedua bentuk tersebut ditampilkan pada
Kode~\ref{lst:bongkar-ukuran}. Karena program berbentuk \file{.COM},
seluruh instruksi dimulai pada offset \hex{100}.

\begin{session}[caption={Hasil perintah \code{U} pada
\file{ukuran-instruksi.com}, hanya bagian kedua bentuk.},
label={lst:bongkar-ukuran}]
-u 100 10A
13F2:0100 B80100        MOV     AX,0001     ; bentuk A, 3 byte
13F2:0103 BB0100        MOV     BX,0001     ; bentuk A, 3 byte
13F2:0106 BB0100        MOV     BX,0001     ; bentuk B, 3 byte
13F2:0109 89D8          MOV     AX,BX       ; bentuk B, 2 byte
\end{session}

Bentuk A menempati alamat \hex{100} sampai \hex{105}, yaitu enam
byte. Bentuk B menempati alamat \hex{106} sampai \hex{10A}, yaitu
lima byte. Keduanya berakhir dengan \reg{AX} = \dec{1} dan
\reg{BX} = \dec{1}, sehingga hasil akhirnya identik. Selisih satu
byte berasal dari \eng{immediate} 16 bit kedua yang tidak lagi
diperlukan. Rincian selisih tersebut diberikan pada
Tabel~\ref{tab:dua-bentuk}.

\begin{table}[H]
    \centering
    \caption{Perbandingan jumlah byte dua bentuk penulisan yang
    setara.}
    \label{tab:dua-bentuk}
    \small
    \begin{tabular}{@{}p{2.0cm}p{4.6cm}p{3.0cm}c@{}}
        \toprule
        \textbf{Bentuk} & \textbf{Instruksi} &
        \textbf{Kode mesin} & \textbf{Byte} \\
        \midrule
        A & \lstasm|MOV AX, 0001h| & \hex{B8 01 00} & 3 \\
        \addlinespace
        A & \lstasm|MOV BX, 0001h| & \hex{BB 01 00} & 3 \\
        \addlinespace
        \multicolumn{3}{@{}l}{\textbf{Jumlah bentuk A}} & \textbf{6} \\
        \midrule
        B & \lstasm|MOV BX, 0001h| & \hex{BB 01 00} & 3 \\
        \addlinespace
        B & \lstasm|MOV AX, BX| & \hex{89 D8} & 2 \\
        \addlinespace
        \multicolumn{3}{@{}l}{\textbf{Jumlah bentuk B}} & \textbf{5} \\
        \bottomrule
    \end{tabular}
\end{table}

\begin{infobox}
Bentuk B belum tentu lebih baik dalam segala hal. Ia lebih hemat satu
byte, tetapi bergantung pada nilai register lain dan lebih sulit
dibaca. Selisih kecil seperti ini baru menjadi berarti pada rutin
yang dipanggil ribuan kali atau pada memori yang sangat terbatas.
\end{infobox}

\section{Praktikum}
\label{sec:praktikum-bab7}

\begin{praktikumbox}
\vspace{1.1em}
\textbf{Judul:} Menulis, meng-assemble, men-link, dan menjalankan
program \file{.EXE} serta \file{.COM}.\\[0.3em]
\textbf{Sub-CPMK:} 2.1 \quad\textbullet\quad
\textbf{Alokasi waktu:} 100 menit \quad\textbullet\quad
\textbf{Bentuk:} individu
\end{praktikumbox}

\subsection{Tujuan Praktikum}

Setelah menyelesaikan praktikum ini, mahasiswa mampu menulis program
assembly, meng-assemble dengan TASM, men-link dengan TLINK,
menjalankan program \file{.EXE} dan \file{.COM}, serta memeriksa
berkas hasil dengan DEBUG.

\subsection{Alat dan Bahan}

\begin{itemize}
    \item Komputer dengan emu8086 terpasang.
    \item DOSBox dengan TASM, TLINK, dan DEBUG.
    \item Berkas \file{hello.asm} dan \file{direktif-data.asm} dari
          \file{code/asm/}.
    \item Lembar jurnal praktikum.
\end{itemize}

\subsection{Langkah Kerja}

\begin{enumerate}
    \item \textbf{Menyiapkan program \file{.EXE}.}
    \begin{enumerate}
        \item Buka DOSBox, lalu masuk ke direktori kerja.
        \item Salin \file{hello.asm} ke direktori kerja.
        \item Rakit dengan \file{tasm hello.asm}.
        \item Link dengan \file{tlink hello.obj}.
        \item Jalankan \file{hello} dan catat keluarannya.
    \end{enumerate}

    \item \textbf{Memeriksa berkas hasil.}
    \begin{enumerate}
        \item Catat ukuran \file{hello.obj} dan \file{hello.exe}.
        \item Jelaskan mengapa berkas \file{.OBJ} belum dapat
              dijalankan.
        \item Jelaskan peran TLINK dalam menghasilkan
              \file{hello.exe}.
    \end{enumerate}

    \item \textbf{Membuat program \file{.COM}.}
    \begin{enumerate}
        \item Tulis program berikut sebagai \file{hello-com.asm}.

\begin{lstlisting}[language=MASM, numbers=none]
.MODEL TINY
.CODE
ORG 100h
START:
    JMP MULAI
pesan DB 'Halo, COM!', 0Dh, 0Ah, '$'
MULAI:
    MOV AH, 09h
    MOV DX, OFFSET pesan
    INT 21h
    INT 20h
END START
\end{lstlisting}

        \item Rakit dengan \file{tasm hello-com.asm}.
        \item Link dengan \file{tlink /t hello-com.obj}.
        \item Jalankan \file{hello-com} dan catat keluarannya.
        \item Bandingkan ukuran \file{.COM} dengan \file{.EXE}.
    \end{enumerate}

    \item \textbf{Memeriksa program dengan DEBUG.}
    \begin{enumerate}
        \item Jalankan \file{debug hello.exe}.
        \item Ketik \file{r} dan catat nilai \reg{CS}, \reg{DS},
              dan \reg{IP}.
        \item Ketik \file{d ds:0000 L 20} untuk melihat isi segmen
              data.
        \item Temukan string \code{"Halo, dunia!"} pada hasil
              pemeriksaan.
        \item Keluar dari DEBUG dengan \file{q}.
    \end{enumerate}

    \item \textbf{Menguji direktif data.}
    \begin{enumerate}
        \item Rakit dan link \file{direktif-data.asm}.
        \item Jalankan program dan catat keluarannya.
        \item Buka program dengan DEBUG, lalu periksa isi segmen
              data dengan \file{d ds:0000 L 40}.
        \item Tandai byte yang berasal dari \dir{DB}, \dir{DW},
              \dir{DD}, dan \dir{DUP}.
    \end{enumerate}

    \item \textbf{Memodifikasi program.}
    \begin{enumerate}
        \item Ubah isi \code{pesan} pada \file{hello.asm} menjadi
              nama lengkap Anda.
        \item Ubah \code{bintang DB 4 DUP('*')} menjadi
              \code{6 DUP('*')}.
        \item Rakit ulang kedua program dan pastikan keluarannya
              berubah.
    \end{enumerate}
\end{enumerate}

\subsection{Hasil yang Dikumpulkan}

\begin{enumerate}
    \item Berkas \file{hello.asm}, \file{hello-com.asm}, dan
          \file{direktif-data.asm} versi modifikasi.
    \item Tangkapan layar keluaran ketiga program.
    \item Tabel perbandingan ukuran \file{.OBJ}, \file{.EXE}, dan
          \file{.COM}.
    \item Tangkapan layar pemeriksaan segmen data dengan DEBUG.
    \item Penjelasan singkat perbedaan program \file{.COM} dan
          \file{.EXE}.
\end{enumerate}

\begin{asesmenbox}
\vspace{1.1em}
\textbf{Asesmen praktikum 5 (bobot 5\% dari komponen praktikum).}

\begin{tabular}{@{}p{6.2cm}cp{4.6cm}@{}}
    \toprule
    \textbf{Indikator} & \textbf{Bobot} & \textbf{Kriteria} \\
    \midrule
    Program \file{.EXE} berhasil dirakit dan dijalankan & 30\% &
    Berkas \file{.OBJ} dan \file{.EXE} terbentuk, keluaran benar \\
    \addlinespace
    Program \file{.COM} berhasil dibuat & 25\% &
    Berkas \file{.COM} terbentuk dan berjalan dengan benar \\
    \addlinespace
    Pemeriksaan segmen data dengan DEBUG & 25\% &
    String dan data ditemukan serta ditafsirkan dengan benar \\
    \addlinespace
    Analisis perbedaan COM dan EXE & 20\% &
    Perbedaan segmen, ukuran, dan inisialisasi DS dijelaskan \\
    \bottomrule
\end{tabular}
\end{asesmenbox}

\section{Latihan}
\label{sec:latihan-bab7}

\begin{enumerate}
    \item \textbf{Konseptual.} Jelaskan hubungan antara
          \eng{mnemonic} dan \eng{opcode}. Mengapa programmer cukup
          menulis \lstasm|MOV AX, BX| tanpa menentukan sendiri byte
          \hex{89 D8}?

    \item \textbf{Konseptual.} Sebutkan tiga sebab yang membuat satu
          \eng{mnemonic} seperti \asm{MOV} memiliki banyak bentuk
          \eng{encoding}. Berikan satu contoh instruksi untuk
          masing-masing sebab.

    \item \textbf{Analisis field.} Sebuah instruksi memiliki byte
          \hex{C7 80 34 12 78 56}. Tentukan:
          \begin{enumerate}
              \item \eng{opcode}-nya dan operasi yang dilakukan;
              \item nilai biner field \code{mod}, \code{reg}, dan
                    \code{r/m};
              \item panjang dan isi \eng{displacement};
              \item panjang dan isi \eng{immediate}.
          \end{enumerate}

    \item \textbf{Encoding manual.} Tuliskan byte kode mesin untuk
          instruksi berikut beserta langkah perhitungannya.
          \begin{enumerate}
              \item \lstasm|MOV DX, 00FFh|
              \item \lstasm|ADD AL, 3|
              \item \lstasm|MOV [BX], AL|
          \end{enumerate}

    \item \textbf{Analisis ModR/M.} Uraikan byte ModR/M berikut
          menjadi ketiga fieldnya, lalu jelaskan operand apa yang
          diwakilinya.
          \begin{enumerate}
              \item \hex{D8}
              \item \hex{44}
              \item \hex{C1}
              \item \hex{16}
          \end{enumerate}

    \item \textbf{Displacement dan immediate.} Jelaskan perbedaan
          peran \eng{displacement} dan \eng{immediate}. Kapan
          keduanya muncul pada sebuah instruksi, dan dapatkah
          keduanya muncul sekaligus? Berikan contoh.

    \item \textbf{Analisis ukuran dan kecepatan.} Dua program memiliki
          hasil akhir sama. Program pertama memakai
          \lstasm|MOV AX, 0001h| dan \lstasm|MOV BX, 0001h|; program
          kedua memakai \lstasm|MOV BX, 0001h| dan \lstasm|MOV AX, BX|.
          \begin{enumerate}
              \item Hitung jumlah byte kedua program.
              \item Sebutkan field yang menjadi penyebab selisihnya.
              \item Kaitkan dengan jumlah byte yang harus diambil
                    prosesor.
          \end{enumerate}

    \item \textbf{Analisis kode dan kesalahan.} Diberikan keluaran
          perintah \code{U} berikut.

\begin{session}[numbers=none]
1A2B:0100 B80500        MOV     AX,0005
1A2B:0103 8908          MOV     [BX+SI],CX
1A2B:0105 CD21          INT     21
\end{session}

          \begin{enumerate}
              \item Jelaskan mengapa baris kedua dapat berakhir pada
                    alamat \hex{103}, bukan \hex{102}.
              \item Mengapa disassembler menampilkan
                    \lstasm|[BX+SI]| dan bukan nama variabel yang
                    mungkin dipakai pada kode sumber?
              \item Sebutkan satu risiko bila hasil \eng{disassembly}
                    langsung dipercaya tanpa memeriksa kode
                    sumbernya.
          \end{enumerate}
\end{enumerate}

\begin{refleksibox}
Jawab pertanyaan berikut pada jurnal praktikum Anda.
\begin{enumerate}
    \item Bagian mana yang paling menantang saat Anda mengisi lembar
          kerja \eng{encoding} secara manual: menentukan
          \eng{opcode}, menyusun ModR/M, atau mengurutkan
          \eng{little-endian}? Jelaskan alasannya.
    \item Setelah melihat hasil perintah \code{U}, apakah dugaan awal
          Anda tentang panjang sebuah instruksi berubah? Bagian mana
          yang paling mengejutkan?
    \item Konsep apa pada bab ini yang masih perlu Anda latih lagi
          sebelum mempelajari mode pengalamatan pada bab berikutnya?
\end{enumerate}
\end{refleksibox}

\section{Rangkuman}
\label{sec:rangkuman-bab7}

Pada bab ini telah dibahas:

\begin{itemize}
    \item Program assembly melewati rantai proses: berkas
          \file{.ASM} ditulis di editor, dirakit oleh TASM menjadi
          \file{.OBJ}, di-link oleh TLINK menjadi \file{.EXE} atau
          \file{.COM}, lalu dimuat DOS ke memori.

    \item TASM menerjemahkan instruksi dan direktif menjadi berkas
          objek. TLINK menyelesaikan rujukan simbol dan menetapkan
          alamat segmen.

    \item Program \file{.COM} memakai satu segmen, berukuran
          maksimum sekitar 64 KB, memakai \dir{ORG 100h}, dan tidak
          memerlukan inisialisasi \reg{DS}. Program \file{.EXE}
          dapat memakai beberapa segmen dan memerlukan inisialisasi
          \reg{DS}.

    \item Direktif \dir{.MODEL} menentukan model memori. Model
          \code{TINY} dipakai untuk program \file{.COM}, sedangkan
          \code{SMALL} dipakai untuk program \file{.EXE} sederhana.

    \item Kerangka program memuat \dir{.MODEL}, \dir{.STACK},
          \dir{.DATA}, \dir{.CODE}, label titik masuk, dan
          \dir{END}.

    \item \dir{ORG 100h} menetapkan awal program \file{.COM} pada
          offset \hex{100} karena 256 byte pertama dipakai
          \eng{Program Segment Prefix}.

    \item Label menandai lokasi instruksi atau data. Komentar
          ditulis setelah tanda \code{;} dan diabaikan assembler.

    \item Direktif data \dir{DB}, \dir{DW}, dan \dir{DD} menentukan
          ukuran data. \dir{EQU} memberi nama konstanta, sedangkan
          \dir{DUP} mengulang pola data.

    \item Kesalahan penulisan dilaporkan TASM beserta nomor baris.
          Kesalahan logika tidak selalu dilaporkan dan harus dicari
          dengan \eng{debugger}.
\end{itemize}

\begin{checklistbox}
\vspace{1.1em}
Tandai setiap butir yang sudah Anda kuasai.

\begin{itemize}
    \item[\checkbox] Saya dapat menjelaskan rantai proses
          \file{.ASM} $\rightarrow$ \file{.OBJ} $\rightarrow$
          \file{.EXE}/\file{.COM}.
    \item[\checkbox] Saya dapat menjelaskan peran TASM dan TLINK.
    \item[\checkbox] Saya dapat menyebutkan perbedaan program
          \file{.COM} dan \file{.EXE}.
    \item[\checkbox] Saya dapat menjelaskan fungsi \dir{.MODEL},
          \dir{.STACK}, \dir{.DATA}, \dir{.CODE}, dan \dir{END}.
    \item[\checkbox] Saya dapat menjelaskan peran \dir{ORG 100h}
          pada program \file{.COM}.
    \item[\checkbox] Saya dapat menulis label dan komentar dengan
          benar.
    \item[\checkbox] Saya dapat menggunakan \dir{DB}, \dir{DW},
          \dir{DD}, \dir{EQU}, dan \dir{DUP}.
    \item[\checkbox] Saya dapat meng-assemble, men-link, dan
          menjalankan program sederhana.
\end{itemize}

\vspace{0.4em}
Bila masih ada butir yang belum bertanda, ulangi bagian yang
bersesuaian sebelum melanjutkan ke bab berikutnya.
\end{checklistbox}

\begin{infobox}
\vspace{1.1em}
\textbf{Rujukan bab.} Materi pada bab ini diringkas dari
\cite{gunawan2019}, Bagian 4, hlm.~26--34, dan Bagian 8,
hlm.~76--77, serta \cite{sto2001}, BAB V--VI, hlm.~21--26, dan
BAB IX, hlm.~41--46.
\end{infobox}

\section{Evaluasi}
\label{sec:evaluasi-bab7}

\subsection{Soal Pilihan Ganda}

\begin{enumerate}
    \item Berkas keluaran yang dihasilkan langsung oleh TASM adalah
          \dots
          \begin{enumerate}
              \item \file{.ASM} \item \file{.OBJ}
              \item \file{.EXE} \item \file{.COM}
          \end{enumerate}

    \item Perangkat yang menggabungkan berkas objek menjadi berkas
          yang dapat dijalankan adalah \dots
          \begin{enumerate}
              \item editor \item TASM
              \item TLINK \item DEBUG
          \end{enumerate}

    \item Direktif yang wajib dipakai pada program \file{.COM}
          adalah \dots
          \begin{enumerate}
              \item \dir{.MODEL SMALL} \item \dir{ORG 100h}
              \item \dir{.STACK 100h} \item \dir{ASSUME CS:CODE}
          \end{enumerate}

    \item Direktif yang menyediakan ruang dua byte untuk satu nilai
          adalah \dots
          \begin{enumerate}
              \item \dir{DB} \item \dir{DW}
              \item \dir{DD} \item \dir{EQU}
          \end{enumerate}

    \item Inisialisasi \reg{DS} diperlukan pada program \file{.EXE}
          karena \dots
          \begin{enumerate}
              \item DOS tidak menetapkan \reg{DS} ke segmen data
                    program
              \item \reg{DS} hanya boleh diisi oleh TLINK
              \item program \file{.EXE} tidak memiliki segmen data
              \item \reg{DS} harus sama dengan \reg{CS}
          \end{enumerate}
\end{enumerate}

\subsection{Soal Uraian}

\begin{enumerate}
    \item Jelaskan rantai proses dari berkas \file{.ASM} sampai
          program berjalan. Sebutkan peran editor, TASM, TLINK, dan
          \eng{loader}.

    \item Sebutkan empat perbedaan program \file{.COM} dan
          \file{.EXE}. Jelaskan mengapa program \file{.COM} tidak
          memerlukan inisialisasi \reg{DS}, sedangkan program
          \file{.EXE} memerlukannya.

    \item Jelaskan fungsi \dir{DB}, \dir{DW}, \dir{DD}, \dir{EQU},
          dan \dir{DUP}. Berikan satu contoh pemakaian untuk
          masing-masing direktif.

    \item Perhatikan program berikut.
          \begin{lstlisting}[language=MASM, numbers=none]
.MODEL SMALL
.STACK 100h
.DATA
    nilai DW 1234h
.CODE
START:
    MOV AX, @DATA
    MOV DS, AX
    MOV AX, nilai
    MOV AH, 4Ch
    INT 21h
END START
          \end{lstlisting}
          Jelaskan fungsi setiap direktif dan label pada program
          tersebut. Jelaskan pula apa yang terjadi bila
          \lstasm|MOV DS, AX| dihapus.
\end{enumerate}

\subsection{Kisi-kisi Penilaian}

\begin{table}[H]
    \centering
    \caption{Kisi-kisi evaluasi Bab~\ref{chap:memulai-assembly}.}
    \label{tab:kisi-bab7}
    \small
    \begin{tabular}{@{}p{1.5cm}p{5.0cm}p{3.0cm}c@{}}
        \toprule
        \textbf{Sub-CPMK} & \textbf{Indikator yang diukur} &
        \textbf{Bentuk soal} & \textbf{Bobot} \\
        \midrule
        2.1 & Menjelaskan rantai proses dan peran TASM/TLINK &
        Pilihan ganda 1, 2; Uraian 1 & 30\% \\
        \addlinespace
        2.1 & Membedakan program COM dan EXE &
        Pilihan ganda 3, 5; Uraian 2 & 30\% \\
        \addlinespace
        2.1 & Menjelaskan struktur program dan direktif &
        Pilihan ganda 4; Uraian 3, 4 & 40\% \\
        \bottomrule
    \end{tabular}
\end{table}

\begin{checklistbox}
\vspace{1.1em}
\textbf{Checklist Pencapaian Kompetensi
Bab~\ref{chap:memulai-assembly}}

Seorang mahasiswa dinyatakan kompeten pada Sub-CPMK 2.1 apabila:
\begin{itemize}
    \item[\checked] dapat menjelaskan alur
          \file{.ASM} $\rightarrow$ \file{.OBJ} $\rightarrow$
          \file{.EXE}/\file{.COM};
    \item[\checked] dapat menjelaskan peran TASM dan TLINK;
    \item[\checked] dapat membedakan program \file{.COM} dan
          \file{.EXE};
    \item[\checked] dapat menjelaskan fungsi \dir{.MODEL},
          \dir{.STACK}, \dir{.DATA}, \dir{.CODE}, \dir{ORG 100h},
          dan \dir{END};
    \item[\checked] dapat menulis label dan komentar dengan benar;
    \item[\checked] dapat menggunakan direktif data \dir{DB},
          \dir{DW}, \dir{DD}, \dir{EQU}, dan \dir{DUP};
    \item[\checked] dapat meng-assemble, men-link, dan menjalankan
          program sederhana.
\end{itemize}
\end{checklistbox}


\end{document}
